\documentclass[12pt]{article}
\usepackage{amsmath}
\usepackage{amssymb}
\usepackage{geometry}
\usepackage{xcolor}
\usepackage{enumitem}
\usepackage{tcolorbox}

\geometry{margin=1in}

\title{Problem 5: Thinning Property and Independence}
\author{Detailed Explanation}
\date{}

\begin{document}

\maketitle

\section*{The Key Question: When Does Thinning Give Independence?}

The \textbf{thinning property} is crucial for understanding Problem 5. Let's clarify when it applies.

\begin{tcolorbox}[colback=blue!5!white,colframe=blue!75!black,title=Thinning Property]
If $\{N_t\}$ is a Poisson process with rate $\lambda$, and each arrival is independently classified as ``Type A'' with probability $p$ or ``Type B'' with probability $1-p$, then:
\begin{itemize}
    \item Type A arrivals form a Poisson process with rate $\lambda p$
    \item Type B arrivals form a Poisson process with rate $\lambda(1-p)$
    \item The two processes are \textbf{independent}
\end{itemize}
\end{tcolorbox}

\section*{Problem 5a: Fixed $n$ (NO Thinning, NOT Independent)}

\textbf{Setup:}
\begin{itemize}
    \item $n$ button presses total (FIXED)
    \item Each press: Red with prob $1/3$, Green with prob $2/3$
    \item $R_n =$ number of red presses, $G_n =$ number of green presses
\end{itemize}

\subsection*{Why $R_n$ and $G_n$ are NOT Independent}

\textbf{Constraint:} $R_n + G_n = n$ \textbf{always holds}

This is a \textbf{deterministic constraint}:
\begin{itemize}
    \item If $R_n = 5$ and $n = 10$, then $G_n = 5$ with probability 1
    \item Knowing $R_n$ tells you $G_n$ exactly
    \item Perfect negative correlation
\end{itemize}

\textbf{Mathematical proof of dependence:}

For independence, we need:
\begin{equation*}
P(R_n = r, G_n = g) = P(R_n = r) \cdot P(G_n = g) \quad \text{for all } r, g
\end{equation*}

But this \textbf{fails}. Consider $n = 10$, $r = 3$, $g = 4$:

\textbf{Left side:}
\begin{equation*}
P(R_{10} = 3, G_{10} = 4) = 0 \quad \text{(since $3 + 4 \neq 10$)}
\end{equation*}

\textbf{Right side:}
\begin{align*}
P(R_{10} = 3) &= \binom{10}{3}\left(\frac{1}{3}\right)^3\left(\frac{2}{3}\right)^7 > 0\\
P(G_{10} = 4) &= \binom{10}{4}\left(\frac{2}{3}\right)^4\left(\frac{1}{3}\right)^6 > 0\\
P(R_{10} = 3) \cdot P(G_{10} = 4) &> 0
\end{align*}

So $0 \neq P(R_{10} = 3) \cdot P(G_{10} = 4)$, violating independence.

\begin{tcolorbox}[colback=red!5!white,colframe=red!75!black,title=Answer for 5a]
\textbf{Joint Distribution:}
\begin{equation*}
P(R_n = r, G_n = g) = \begin{cases}
\binom{n}{r}\left(\frac{1}{3}\right)^r\left(\frac{2}{3}\right)^{n-r} & \text{if } r + g = n\\
0 & \text{otherwise}
\end{cases}
\end{equation*}

\textbf{Independence:} \textcolor{red}{\textbf{NO}}, $R_n$ and $G_n$ are \textbf{NOT independent}
\end{tcolorbox}

\section*{Why Doesn't Thinning Apply Here?}

Thinning applies to \textbf{Poisson processes} (continuous time, random number of arrivals).

In Problem 5a:
\begin{itemize}
    \item We have a \textbf{fixed} number $n$ of trials
    \item This is \textbf{binomial}, not Poisson
    \item No thinning property
\end{itemize}

Think of it this way: With $n = 10$ button presses, once you count the red presses, the number of green presses is determined. There's no randomness left!

\newpage

\section*{Problem 5b: Random $N$ (Thinning DOES Apply, ARE Independent)}

\textbf{Setup:}
\begin{itemize}
    \item $N \sim \text{Poisson}(a)$ total presses (RANDOM)
    \item Each press: Red with prob $1/3$, Green with prob $2/3$
    \item $R_N =$ number of red presses, $G_N =$ number of green presses
\end{itemize}

\subsection*{Why $R_N$ and $G_N$ ARE Independent}

Now we \textbf{can} apply the thinning property!

Think of this as a Poisson process:
\begin{itemize}
    \item Start with a Poisson process of button presses with rate $\lambda$ over time $[0, T]$ where $\lambda T = a$
    \item Each arrival is independently colored red (prob $1/3$) or green (prob $2/3$)
    \item By the \textbf{thinning property}:
    \begin{itemize}
        \item Red presses form a Poisson process with rate $\lambda \cdot \frac{1}{3}$
        \item Green presses form a Poisson process with rate $\lambda \cdot \frac{2}{3}$
        \item These two processes are \textbf{independent}
    \end{itemize}
\end{itemize}

Over the interval, the total counts are:
\begin{align*}
R_N &\sim \text{Poisson}\left(\frac{a}{3}\right)\\
G_N &\sim \text{Poisson}\left(\frac{2a}{3}\right)
\end{align*}

And they are \textbf{independent}!

\begin{tcolorbox}[colback=green!5!white,colframe=green!75!black,title=Answer for 5b]
\textbf{Joint Distribution:}
\begin{equation*}
P(R_N = r, G_N = g) = \frac{(a/3)^r e^{-a/3}}{r!} \cdot \frac{(2a/3)^g e^{-2a/3}}{g!}
\end{equation*}

\textbf{Independence:} \textcolor{green!50!black}{\textbf{YES}}, $R_N$ and $G_N$ \textbf{ARE independent}
\end{tcolorbox}

\subsection*{Why Is This Different?}

The key difference:
\begin{itemize}
    \item In 5a: $n$ is \textbf{fixed}, so $R_n + G_n = n$ is a rigid constraint
    \item In 5b: $N$ is \textbf{random}, so $R_N + G_N = N$ is itself random
    \item The randomness of $N$ breaks the dependency!
\end{itemize}

\textbf{Intuition:}
\begin{itemize}
    \item With fixed $n$: Once you know how many reds, you know how many greens (no randomness left)
    \item With random $N$: Even if you know $R_N = 5$, you don't know $G_N$ because you don't know $N$
\end{itemize}

\section*{Summary Table}

\begin{center}
\begin{tabular}{|l|c|c|c|}
\hline
\textbf{Problem} & \textbf{Total} & \textbf{Thinning?} & \textbf{Independent?}\\
\hline
5a: $R_n$, $G_n$ & Fixed $n$ & No (Binomial) & \textcolor{red}{\textbf{NO}}\\
\hline
5b: $R_N$, $G_N$ & Random $N \sim \text{Poisson}(a)$ & Yes (Poisson) & \textcolor{green!50!black}{\textbf{YES}}\\
\hline
\end{tabular}
\end{center}

\vspace{1em}

\begin{tcolorbox}[colback=yellow!5!white,colframe=orange!75!black,title=Key Takeaway]
\textbf{Thinning property} gives independence when:
\begin{itemize}
    \item You start with a \textbf{Poisson process}
    \item You randomly classify arrivals
\end{itemize}

It does \textbf{NOT} apply when:
\begin{itemize}
    \item You have a \textbf{fixed} number of trials (binomial setting)
    \item The constraint $R_n + G_n = n$ makes them dependent
\end{itemize}
\end{tcolorbox}

\end{document}
